\subsection{问题一分析}
问题一的核心不仅在于验证"鱼类资源量增加"，更在于同时揭示"鱼类增长的驱动机制"与"基础资源的承载压力变化"。若将四大家鱼合并为单一消费者功能群，将水草、浮游植物、浮游动物和底栖饵料简化为单一资源总量，虽仍能匹配2025年资源量恢复至1.8倍的观测事实，但无法揭示哪一层级资源率先承压，也难以解释题目中"水下荒漠"等底层生态退化现象。因此，模型必须保留四类基础资源与四大家鱼的基本食性对应关系。

题目给出的两类观测数据需差异化使用：2025年四大家鱼1.8倍恢复量对应全流域平均口径，可用于约束系统总量参数；湖北江段单网次渔获量提升120%具有明显局地性和网具特异性，更适合作为外部独立验证，不宜与全流域总量数据共同参与参数标定。据此，我们先在统一总量尺度下识别禁渔的直接生态效应，再将局地捕获差异纳入结果分析。

